\documentclass[11pt]{amsart}
\usepackage{calc,amssymb,amsmath,ulem,stmaryrd,fullpage}
\usepackage{enumerate}
\usepackage{alltt}
\RequirePackage[dvipsnames,usenames]{color}
\let\mffont=\sf
\normalem
%\input{kmacros3.sty}
\input{xy}
\xyoption{all}
\usepackage{hyperref}
\usepackage{graphicx}
\usepackage[all,cmtip]{xy}
\usepackage{verbatim}
\usepackage[left=0.95in,top=0.95in,right=0.95in,bottom=0.95in]{geometry}
\newcommand{\tauCohomology}{T}
\newcommand{\FCohomology}{S}
\usepackage{mathrsfs}


\theoremstyle{theorem}
%\newtheorem*{mainthm*}{Main Theorem}

\newcommand{\note}[1]{\marginpar{\sffamily\tiny #1}}

\def\todo#1{\textcolor{Mahogany}%
{\sffamily\footnotesize\newline{\color{Mahogany}\fbox{\parbox{\textwidth-15pt}{#1}}}\newline}}

\renewcommand{\O}{\mathcal O}
\newcommand{\ie}{{\itshape i.e.} }
\newcommand{\roundup}[1]{\lceil #1 \rceil}
\newcommand{\rounddown}[1]{\lfloor #1 \rfloor}
\renewcommand{\to}[1][]{\xrightarrow{\ #1\ }}
\newcommand{\onto}[1][]{\protect{\xrightarrow{\ #1\ }\hspace{-0.8em}\rightarrow}}
\newcommand{\into}[1][]{\lhook \joinrel \xrightarrow{\ #1\ }}
%\newcommand{DBDual}[1]{{\underline{\omega}_{#1}}}
\newcommand{\DBDual}[1]{{{\uuline \omega}{}^{\mydot}_{#1}}}
\newcommand{\Tt}{{\mathfrak{T}}}
%\newcommand{\bn}{{\mathfrak{n}} }
\newcommand{\sol}[1]{ \vskip 6pt{\color{blue} {\bf Solution:} #1} \vskip 6pt}

\newcommand{\bR}{{\mathbb{R}}}
\newcommand{\va}{{\vec{a}}}
\newcommand{\vb}{{\vec{b}}}
\newcommand{\vzero}{{\vec{0}}}
\newcommand{\vi}{{\vec{i}}}
\newcommand{\vj}{{\vec{j}}}
\newcommand{\vk}{{\vec{k}}}
\newcommand{\vh}{{\vec{h}}}
\newcommand{\vr}{{\vec{r}}}
\newcommand{\vu}{{\vec{u}}}
\newcommand{\vv}{{\vec{v}}}
\newcommand{\vw}{{\vec{w}}}
\newcommand{\vx}{{\vec{x}}}
\newcommand{\vy}{{\vec{y}}}
\newcommand{\vz}{{\vec{z}}}
\newcommand{\vT}{{\vec{T}}}
\newcommand{\vN}{{\vec{N}}}
\newcommand{\vF}{{\vec{F}}}
\newcommand{\vG}{{\vec{G}}}
\newcommand{\bF}{\mathbb{F}}
\newcommand{\bQ}{\mathbb{Q}}
\newcommand{\bZ}{\mathbb{Z}}
\newcommand{\bA}{\mathbb{A}}
%\newcommand{\cF}{\mathcal{F}}
%\newcommand{\cG}{\mathcal{G}}
%\newcommand{\cH}{\mathcal{H}}
\newcommand{\image}{\mathrm{im}\,}
\newcommand{\coker}{\mathrm{coker}\,}
\newcommand{\Hom}{\mathrm{Hom}\,}
\newcommand{\Supp}{\mathrm{Supp}\,}
\newcommand{\Ann}{\mathrm{Ann}\,}
\newcommand{\sH}{\mathscr{H}\,}
\newcommand{\cI}{\mathcal{I}}
\newcommand{\cJ}{\mathcal{J}}
\newcommand{\cE}{\mathcal{E}}
\newcommand{\cF}{\mathcal{F}}
\newcommand{\cG}{\mathcal{G}}
\newcommand{\cH}{\mathcal{H}}
\newcommand{\cO}{\mathcal{O}}
\newcommand{\cM}{\mathcal{M}}
\newcommand{\cN}{\mathcal{N}}

\begin{document}
\title{Worksheet \#2 -- Math 6140\\Spring 2019}
\author{Due Friday, February 1st (electronically)}

\maketitle
You may work in groups of up to 3.  Only one worksheet needs to be turned in per group.
\vskip 9pt
\noindent
{\bf 1.}  Suppose that $\phi : Y \to X$ is a regular map of quasi-projective varieties and that $\cE$ a locally free sheaf of $\O_X$-modules.  Prove that $\phi^* \cE$ is also locally free as a sheaf of $\O_Y$-modules.
\vskip 200pt
\noindent
{\bf 2.}  Suppose that $\phi : Y \to X$ is a regular map of quasi-projective varieties.  Suppose that $\cE$ is a locally free sheaf of $\O_X$-modules on $X$.  Further suppose that $\cF$ is a sheaf of $\O_Y$-modules.  Prove that there is a natural isomorphism of $\O_Y$-modules
\[
(\phi_* \cF) \otimes_{\cO_X} \cE \cong \phi_* (\cF  \otimes_{\cO_Y} \phi^* \cE).
\]
(don't actually prove the ``natural'' part, but at least state what you would have to show)
\newpage
\noindent
{\bf 3.}  Let $X$ be a quasi-projective variety and suppose that $\cF$ is a coherent sheaf on $X$.  Suppose that there is another coherent sheaf $\cG$ on $X$ such that
\[
\cF \otimes_{\cO_X} \cG \cong \cO_X.
\]
Prove that $\cF$ is locally free and of rank 1.
\vskip 3pt
\emph{Hint:} Nakayama's lemma and the ideas surrounding it may help.
\vskip 250pt
\noindent
{\bf 4.}  Suppose that $X$ is a quasi-projective variety and that $\cF$ is a locally free sheaf of rank $1$ on $X$.  Prove that there exists a coherent sheaf $\cG$ so that
\[
\cF \otimes_{\cO_X} \cG \cong \cO_X.
\]
In view of this problem and its predecessor, locally free sheaves of rank 1 are called \emph{invertible sheaves}.
\newpage
\noindent
{\bf 5.} Suppose $X$ is an affine variety and that $R = k[X]$.  Suppose that $M$ is an $R$-module.  For any $m \in M = \Gamma(X, \widetilde{M})$, show that $\Supp m = Z(\Ann_R m)$.  Here $\Ann_R m = 0 :_R m = \{r \in R \;|\; rm = 0 \}$ and $\Supp m$ is defined as in the previous worksheet.
\vskip 250pt
\noindent
{\bf 6.}  With notation as above, if $M$ is finitely generated, show that $\Supp \widetilde M = Z(\Ann_R M)$.
\end{document}

